\documentclass[conference]{IEEEtran}

\usepackage{multirow}
\usepackage{float}
\usepackage{cite}
\usepackage{caption}
\usepackage{subcaption}
\usepackage{amsmath,amssymb,amsfonts,amstext}
% \usepackage{algorithmic}
\usepackage[dvipdfmx]{graphicx}
\usepackage{textcomp}
\usepackage{xcolor}
\usepackage{colortbl}
\usepackage[T1]{fontenc}
\usepackage{mdframed}
\usepackage{algorithm, algpseudocode} % texlive-science
%\usepackage[braket, qm]{qcircuit}
\usepackage{newtxtext}
\usepackage[normalem]{ulem}
\usepackage{nameref}
\usepackage{tikz}
\usetikzlibrary{arrows.meta, shapes, positioning, calc, fit,tikzmark,quantikz2}
\usepackage{quantikz}

\newcommand*\rtofstarg{\tikz[baseline=(char.base)]{
            \node[shape=circle,double,draw,inner sep=1pt,minimum width=16pt] (char) {S};}}
\newcommand*\rtofsdgtarg{\tikz[baseline=(char.base)]{
            \node[shape=circle,double,draw,inner sep=1pt] (char) {S$^{\dagger}$};}}
\newcommand*\rtoftarg{\tikz[baseline=(char.base)]{
            \node[shape=circle,draw,inner sep=-0.5pt,fill=none] (char) {$\oplus$};}}
\newcommand*\rtofsctrl{\tikz[baseline=(char.base)]{
            \node[shape=circle,double,draw,inner sep=1pt,fill=black, minimum width=3pt] (char) {};}}
\newcommand*\rtofsctrlo{\tikz[baseline=(char.base)]{
            \node[shape=circle,double,draw,inner sep=1pt,fill=white, minimum width=3pt] (char) {};}}
%\newcommand*\rtofsctrlo{\tikz[baseline=(char.base)]{
%            \node[shape=circle,double,draw,inner sep=1pt,fill=white, minimum width=3pt] (char) {};}}
\newcommand*\dashedwire{\tikz{
            \draw[dashed] (0,0) -- (0.5,0); }}

\newcommand*\ctrlrtofs[1]{\push{\rtofsctrl}\qwx[#1]\qw}
\newcommand*\ctrlortofs[1]{\push{\rtofsctrlo}\qwx[#1]\qw}
\newcommand*\targrtofs{\push{\rtofstarg}\qw}
\newcommand*\targrtof{\push{\rtoftarg}\qw}
% \newcommand*\qcbox[1,2,3,4]{\gategroup[4,steps=2,style={dashed,rounded corners,fill=blue!20, inner xsep=2pt},background,label style={label position=below,anchor=north,yshift=-0.2cm}]{Compute}}


\makeatletter
\newcommand{\sublabel}[1]{\protected@edef\@currentlabel{\thefigure(\thesubfigure)}\label{#1}}
\makeatother
\captionsetup{subrefformat=parens}
\algrenewcommand\algorithmicindent{0.5em}
\newtheorem{definition}{Definition}
\newtheorem{example}{Example}
\newcommand{\thickhline}{%
    \noalign {\ifnum 0=`}\fi \hrule height 1pt
    \futurelet \reserved@a \@xhline
}
%\def\BibTeX{{\rm B\kern-.05em{\sc i\kern-.025em b}\kern-.08em
%    T\kern-.1667em\lower.7ex\hbox{E}\kern-.125emX}}
\newcommand{\figcaption}[1]{\def\@captype{figure}\caption{#1}}
\newcommand{\tblcaption}[1]{\def\@captype{table}\caption{#1}}

\newcommand{\Del}[1]{\textcolor{red}{\sout{#1}}}
\newcommand{\Add}[1]{\textcolor{blue}{\emph{#1}}}

\newcount\NGCcolor
\NGCcolor=1
\newcount\NGCcolortwo
\NGCcolortwo=1
%\NGCcolor=0
\ifnum\NGCcolortwo<1
  \newcommand{\blueoutt}[1]{\textcolor{blue}{#1}}  
\else
  \newcommand{\blueoutt}[1]{#1}
\fi
  
\ifnum\NGCcolor<1
  \newcommand{\redout}[1]{\textcolor{red}{#1}}
  \newcommand{\blueout}[1]{\textcolor{blue}{#1}}
  \newcommand{\greenout}[1]{\textcolor{green}{#1}}
  \newcommand{\cyannout}[1]{\textcolor{cyan}{#1}}
  \newcommand{\magentaout}[1]{\textcolor{magenta}{#1}}
  \newcommand{\yellowout}[1]{\textcolor{yellow}{\bf #1}}
  \newcommand{\algblueout}[1]{\leavevmode\color{blue}{#1}}
\else
  \newcommand{\redout}[1]{#1}
  \newcommand{\blueout}[1]{#1}
  \newcommand{\greenout}[1]{#1}
  \newcommand{\cyannout}[1]{#1}
  \newcommand{\magentaout}[1]{#1}
  \newcommand{\yellowout}[1]{#1}
  \newcommand{\algblueout}[1]{\leavevmode {#1}}
\fi

\newcommand{\secref}[1]{Sec.~\ref{#1}}
\newcommand{\figref}[1]{Fig.~\ref{#1}}
\newcommand{\tabref}[1]{Table~\ref{#1}}
\newcommand{\eqnref}[1]{Eq.~\ref{#1}}

\algnewcommand\algorithmicswitch{\textbf{switch}}
\algnewcommand\algorithmiccase{\textbf{case}}
\algnewcommand\algorithmicassert{\texttt{assert}}
\algnewcommand\Assert[1]{\State \algorithmicassert(#1)}%
% New "environments"
\algdef{SE}[SWITCH]{Switch}{EndSwitch}[1]{\algorithmicswitch\ #1\ \algorithmicdo}{\algorithmicend\ \algorithmicswitch}%
\algdef{SE}[CASE]{Case}{EndCase}[1]{\algorithmiccase\ #1}{\algorithmicend\ \algorithmiccase}%
\algtext*{EndSwitch}%
\algtext*{EndCase}%
\renewcommand{\baselinestretch}{1.05}
\begin{document}

\title{XAG Dirty Ancilla Synthesis costs almost the same as Clean Ancilla Synthesis}
\author{
    \IEEEauthorblockN{\phantom{David Clarino}}
    \IEEEauthorblockA{\phantom{Ritsumeikan University} \\
      \phantom{dclarino@fc.ritsumei.ac.jp}}
    \and 
    \IEEEauthorblockN{\phantom{Shigeru Yamashita}}
    \IEEEauthorblockA{\phantom{Ritsumeikan University} \\
      \phantom{ger@cs.ritsumei.ac.jp}}
}

\maketitle

\begin{abstract}
  XOR-AND-Inverter graphs (XAG) have been used to generate quantum Boolean circuits. Using
  techniques such as AND node minimization and AND gate minimization, this expression can
  help reduce the quantum cost to implement such Boolean functions. In particular, in this
  work, we focus on the T-count and the T-depth. However, current XAG generation techniques
  such as caterpillar and LHRS require the use of clean ancilla qubits, or qubits which are
  initialized to the 0 state, which are not as abundant in a real quantum computer.
  This necessitates the use of techniques such as pebbling to conserve the number of clean
  ancillae to be used for synthesis. In this work, we take inspiration from the technique
  of using dirty ancillae, which have some non zero state on them, to decompose Multiple
  Control Toffoli (MCT) gates. In existing constructions, however, dirty ancillae
  decompositions incur additional T-count and T-depth. We demonstrate in this work that
  composing from an XAG incurs mostly the same cost using dirty ancillae as clean ancillae,
  that being a single T-gate. In particular, we demonstrate that their T-counts differ by
  an amount equal to the number of primary outputs. We use this to demonstrate that a
  clean ancilla MCT decomposition and a dirty ancilla MCT decomposition differ in cost by a
  single T-gate, regardless of the number of controls. 
\end{abstract}


%% \begin{abstract}
%%   Multiple Control Toffoli (MCT) gates are an essential part of quantum computing. They create entanglements that
%%   are non-Clifford in nature, enabling quantum computations with advantage over the classical case. MCT gates,
%%   however, are not directly implementable on hardware. They are often decomposed into Toffoli gates, which in
%%   turn are decomposed into the Clifford+T gate set, whose cost can be measured using T-count, the number of
%%   T-gates used to implement a circuit.  With the use of ancilla qubits, MCT gates can be decomposed with a number
%%   of Toffoli gates linear with its inputs, which also means a T-count linear with its inputs. There exist two ways
%%   to do this: clean ancilla, which are initialized to 0, and dirty ancilla, which can have any value. In the
%%   original formulation of the decomposition in Barenco et al, dirty ancilla decompositions use more Toffoli gates
%%   than clean ancilla. Later results, such as Maslov's RTOF decomposition, use relative phase gates, which
%%   implement a Boolean function while allowing a relative phase, to save on T-count. However, even using a relative
%%   phase, these deconstructions show a linear separation in their T-count. Here, we propose a decomposition
%%   which leverages relative phase to further reduce T-count, to the point that it eliminates this linear separation.
%%   That is, we demonstrate that dirty ancilla decompositions of an MCT cost at most 1 more T-gate, regardless of
%%   number of inputs, than clean ancilla decompositions. We demonstrate improvements in T-count for other common
%%   constructions that use this dirty ancilla decomposition, such as the single ancilla decomposition with linear
%%   cost.
%% \end{abstract}

\section{Introduction}
\label{Intro}

Logic network based quantum circuit techniques such as LUT-based Hierarchical Reversible
Synthesis (LHRS)~\cite{bib-lhrs-soeken} and \emph{caterpillar}~\cite{bib-meuli-ros}, which
uses XOR-AND-Inverter graphs (XAG) to represent the Boolean functions. XAG based methods
in particular have a range of techniques to optimize the generated quantum circuit for
T-count, which is the number of T-gates in a quantum circuit, and T-depth, which is the
largest such chain of T-gates in a quantum circuit. These metrics are important to fault
tolerant quantum computing because T-gates have an outsize implementation cost compared
to other gates~\cite{bib-bravyi-kitaev-magic-distillation}. In XAGs, T-depth can be
analogized to AND operations, which are often implemented using other three qubit
non-Clifford gates such as the Toffoli gate and the three qubit controlled-Z gate.
Therefore existing methods to optimize them for AND nodes~\cite{bib-meuli-mult} and
AND-depth\cite{bib-meuli-ros}, which obliquely optimize for T-count and T-depth
respectively.

Existing methods, however, largely use single-target gates acting on clean ancilla,
and are required to be both initialized to $\ket{0}$ and returned to the value at the end
of the calculation. Multiple Control Toffoli (MCT) gates can be decomposed with linear cost
using such ancilla. We can model such a decomposition as an XAG, one of two-input ANDs
chained together. In a similar way, we can also decompose the MCT using \emph{dirty ancilla},
which have no requirements on its initial value. Using dirty ancilla means the circuit
can use any of the other qubits that are presently being used by other subccircuits as
scratch space. This reduces the need for adding clean ancilla towards the total T-count.
In the standard formulation, dirty ancilla decompositions have higher T-count than the
clean ancilla decompositions. Some work optimizes these using a \emph{relative phase},
or a phase on the output that is input dependent, but there remains a linear separation
between the two.

In this work, we derive a way to generate XAG using dirty ancilla, inspired by the dirty
ancilla MCT decomposition. We then use relative phase to further reduce the T-count of
our decomposition, and demonstrate that it is only separated by a T-count equal to the
number of output qubits. This implies that a clean ancilla and dirty ancilla MCT
decomposition only differs by a single T-gate. We conduct several experiments on
EPFL benchmarks to demonstrate this.

This paper is structured as follows. In~\secref{Pre}, we first introduce preliminary
knowledge to understand our proposals, including XAG quantum circuit synthesis and
relative phase. We then present our proposal in~\secref{Pro}, including the dirty
ancilla XAG synthesis, and its implications with MCT decomposition. Finally, \secref{Conc}
concludes with our conclusions and future work.




%% Multiple Control Toffoli (MCT) gates are important components to quantum computing. They
%% allow non-trivial non-Clifford entanglements, which serve as building blocks to quantum
%% algorithms which demonstrate quantum advantage.

%% The canonical Toffoli gate has two control bits and a target bit. It implements the Pauli X
%% operation on its target bit when both of its input qubits are 1. A Toffoli gate is not a
%% gate that can be implemented directly; it is often decomposed into the Clifford+T basis
%% gate set. For fault tolerant computing, the T-gate dominates the implementation cost.
%% Thus, T-count, or the number of T-gates and T$^{\dagger}$-gates in a circuit, is a metric
%% that is often targeted in optimization

%% An MCT gate generalizes the Toffoli gate to multiple control bits. In general, implementing
%% the MCT gate using only the Clifford+T basis gate set without any additional bits takes
%% a T-count that is at best quadratic in the control bits~\cite{bib-barenco-elementary}.
%% However, when ancilla bits are employed, the MCT gates can be decomposed using a linear
%% number of Toffoli gates, and therefore a linear T-count~\cite{bib-barenco-elementary,bib-mike-and-ike}.
%% The methods differ whether the ancilla used are {\it clean ancilla}, which are initialized to $\ket{0}$,
%% or {\it dirty ancilla}, which have no requirements on initial values. Dirty ancilla decompositions
%% are important because they make possible implementations that simply use existing qubits
%% for ancilla, obviating the need to add qubits initialized to $\ket{0}$.

%% Decomposing MCT gates using clean ancilla typically requires a smaller T-count than decomposing
%% using dirty ancilla. In the canonical Barenco et al. formulation, the dirty ancilla decomposition
%% takes almost twice the number of Toffoli as a clean ancilla decomposition, implying a T-count
%% of the same scale. Successive results such as Maslov~\cite{bib-maslov-rtof} brought implied that
%% a T-count less than seven multiplied by the number of Toffoli gates is needed, by using
%% {\it relative phase} constructions. Other results by Abdessaieh et al~\cite{bib-abdessaied-mct}
%% close the gap between the clean ancilla and dirty ancilla decompositions by employing further
%% relative phase style constructions. However, this still has a separation in cost that is
%% linear.

%% In this work, we demonstrate that, by exploiting further relative phase opportunities, clean ancilla
%% and dirty ancilla decompositions have costs that are separated by a constant amount, that being a
%% single T-gate. This decreases the opportunity cost of using the the dirty ancilla decomposition. We
%% use it in this work to improve existing constructions that use the dirty ancilla construction,
%% including the linear cost, single dirty ancilla decomposition by Zindorf and
%% Bose~\cite{bib-zindorf-mct}. We obtain new cost calculations for such decompositions. We summarize
%% the existing state of the art of MCT decompositions in the last chapter.

%% Dctodo
%% This paper is structured as follows: \secref{Pre} goes into the necessary preliminary for this
%% work, including relative phase constructions. \secref{MCT} goes into best known methods of MCT
%% decomposition, many of them using relative phase constructions. \secref{Pro} describes how we
%% derive our proposed dirty ancilla decomposition. \secref{App} shows our application of this
%% proposal to existing methods. \secref{Conc} concludes the paper with our takeaways and possible
%% future research paths.


\begin{figure}[t]
\centering
\setlength{\tabcolsep}{0pt}
\begin{tabular}{
    p{0.48\linewidth}
    @{\hspace{0.01\linewidth}}
    p{0.48\linewidth}
}
  \hspace{-1.0cm}
  \subfloat[X-gate]{%
    \makebox[\linewidth][c]{\begin{quantikz}[row sep=0.2em, column sep=0.5em]
  \lstick{\ket{x}} &\qw & \gate{X}    & \qw & \rstick{\ket{\bar{x}}}  
\end{quantikz}

}%
    \label{fig-xgate}
  }
  &
  \subfloat[Z-gate]{%
    \makebox[\linewidth][c]{\begin{quantikz}[row sep=0.2em, column sep=0.5em]
  \lstick{$\ket{x}$} &\qw & \gate{Z}    & \qw & \rstick{$e^{i \pi x}\ket{x}$}  
\end{quantikz}

}%
    \label{fig-zgate}
  }
\\[0.4em]
\hspace{2.0cm}
\subfloat[H-gate]{%
    \makebox[\linewidth][c]{\begin{quantikz}[row sep=0.2em, column sep=0.5em]
  \lstick{\ket{x}} & \qw & \gate{H}    & \qw & \rstick{$\frac{1}{\sqrt{2}} \sum_{i\in\{0,1\}} (-1)^{i \cdot x}\ket{i}$}
\end{quantikz}

}%
    \label{fig-hgate}
  }
\\[0.4em]
\subfloat[T/T$^{\dagger}$-gate]{%
  \makebox[\linewidth][c]{%
    \vbox{
      \hbox{\input{img/Tgate}}
      \hbox{\begin{quantikz}[row sep=0.2em, column sep=0.5em]
  \lstick{\ket{x}} & \qw & \gate{T^{\dagger}}    & \qw & \rstick{$e^{i\pi \frac{-1}{4} x}\ket{x}$}  
\end{quantikz}

}
    }%
  }
  \label{fig-tgate}
}
&
\subfloat[Toffoli gate]{%
  \makebox[\linewidth][c]{\begin{quantikz}[row sep=0.5em, column sep=0.5em]
  \lstick{\ket{x_0}} & \qw & \ctrl{1} & \qw & \rstick{\ket{x_0}} \\
  \lstick{\ket{x_1}} & \qw & \ctrl{1} & \qw & \rstick{\ket{x_1}} \\
  \lstick{\ket{y}}   & \qw & \targ    & \qw & \rstick{\ket{x_0 \cdot x_1 \oplus y}}  \\
\end{quantikz}

}%
  \label{fig-toff}
}

\\[0.4em]

% Row 4
\subfloat[CNOT-gate]{%
  \makebox[\linewidth][c]{\begin{quantikz}[row sep=0.3cm, column sep=0.3cm]
  \lstick{\ket{x}} & \qw & \ctrl{1} & \qw & \rstick{\ket{x}} \\
  \lstick{\ket{y}} & \qw & \targ    & \qw & \rstick{\ket{x \oplus y}}  \\
\end{quantikz}

}%
  \label{fig-cnot}
}
&
\subfloat[Controlled-Z-gate]{%
  \makebox[\linewidth][c]{\begin{quantikz}[row sep=0.2em, column sep=0.5em]
  \lstick{\ket{x}} & \qw & \ctrl{1}  & \qw & \rstick{\ket{x}} \\
  \lstick{\ket{y}} & \qw & \ctrl{-1} & \qw & \rstick{$(-1)^{x \cdot y}\ket{y}$}  \\
\end{quantikz}

}%
  \label{fig-cz}
}
\\[0.4em]

\end{tabular}

\caption{Quantum gates}
\label{fig-gates}
\end{figure}


\section{Preliminary}
\label{Pre}

\subsection{Quantum Bits and Quantum Gates}
\label{qubits}
Quantum computers internally represent data as \emph{qubits}, which are quantum systems that can
take on the quantum states $\ket{0}$ and $\ket{1}$, each known as \emph{basis states}. The
tensor product of several basis states
$\ket{\bar{x}}=\ket{x_0} \otimes \ket{x_1} \otimes \cdots \otimes \ket{x_{n-1}}, \bar{x} \in \{0,1\}^n$
is known as a \emph{computational basis state}, and a general quantum state can be expressed as a
linear combination of these computational basis states
$\ket{\psi} = k_0\ket{0} + k_1 \ket{1} + k_{n-1} \ket(n-1), k_i = a_ie^{i{\phi}_i} \in \mathbb{C}, \sum_i{k_i * k_i^\dagger} = 1$.


%% These computational basis states can also be in linear superposition.
%% \begin{equation}
%% \ket{\psi} = \sum_{\mathbf{k} \in \{0,1\}^n} \alpha_{\mathbf{k}}\ket{\mathbf{k}},
%% \sum_{\mathbf{k} \in \{0,1\}^n} \alpha_{\mathbf{k}}^2 = 1
%% \end{equation}

{\it Quantum gates} describe qubit mappings $\ket{{\psi}_A} \mapsto \ket{{\psi}_B}$. When a qubit mapping maps
a basis state to another basis state with an input dependent phase
$\ket{\bar{x}} \mapsto \ket{e^{i\phi(\bar{x})}}\ket{\bar{y}}$, $\phi(\bar{x})$ is known as a \emph{relative phase},
and the mapping is said to implement $\ket{\bar{x}} \mapsto \ket{\bar{y}}$ \emph{up to relative phase}. In this
work, we focus on the gate set described in~\figref{fig-gates}.

The CNOT, Z, X, H, and CZ gates are all part of the \emph{Clifford gate set}, which are efficiently simulatable
classically~\cite{bib-gottesman-knill}. Adding $T/T^{\dagger}$ gates creates the \emph{Clifford+T} universal gate
set, which can represent computations that display quantum advantage. However, T-gates have a much higher
implementation cost in fault-tolerant quantum computing~\cite{bib-gottesman-stabilizer}. Thus, {\it T-count}, the
number of $T/T^{\dagger}$ gates, and {\it T-depth}, the longest such serial chain of them in a circuit, are often
used as a cost metric~\cite{amy-meet-in-middle} to measure quantum circuits.

\subsection{Toffoli Gates and Relative Phase Implementations}
\label{sub-toff}

Some gates, such as the Toffoli gate in~\figref{fig-toff}, are composite gates that are implemented using the
Clifford+T gate set. A simple way to derive the Clifford+T implementation of this circuit is to use the symmetry
of the X and Z operations under the Hadamard transformation, and observe that the Toffoli gate is equivalent to
a 3 bit Controlled Z (CCZ) gate, sandwiched between two Hadamards on the qubit that will be the target.

The CCZ gate can then be generated from the CNOT+T gate set, that is, using only CNOTs and T/T$^{\dagger}$ set.
This gate set has a straightforward implementation from what is known as a \emph{sum-of-paths} expression,
generally expressed as below.

\begin{equation}
  \begin{aligned}
    \label{eq-boolean-fourier}
    f(\mathbf{x}) = \sum_{\mathbf{k} \neq 0} \hat{f}(\mathbf{k}) \cdot \chi_k(\mathbf{x}),
    \hat{f}(\mathbf{k}) \in \mathbb{R}
  \end{aligned}
\end{equation}

Where $\chi_k(\mathbf{x}) = ( k_0 x_0 \oplus k_1 x_1 \oplus \cdots \oplus k_{n-1} x_{n-1})$
for $\mathbf{k} \in \{0,1\}^{n}$, or the set of linear Boolean functions over each of the bits
of $\mathbf{x}$, $x_0 \cdots x_{k-1}$.


\begin{figure}[t]
  \centering
  \begin{minipage}{\linewidth}
    \centering
    \scalebox{0.7} {
      %% \Qcircuit @C=0.5em @R=0.2em @!R { \\
%%                              &            &     &          &  & &                           &          &                                &          &                                &          &                             &          & \push{x_a}            &          &     & \\
%%           \lstick{\ket{x_a}} & \ctrl{3}   & \qw &          &  & & \qw                       & \ctrl{4} & \qw                            & \qw      & \qw                            & \ctrl{4} & \qw                         & \qw      & \gate{\mathrm{T}}     & \ctrl{2} & \qw & \\
%%                              &            &     &          &  & &                           &          &                                &          & \push{x_b}                     &          &                             &          & \push{x_a \oplus x_b} &          &     & \\
%%           \lstick{\ket{x_b}} & \ctrl{2}   & \qw & \push{=} &  & & \qw                       & \qw      & \qw                            & \ctrl{2} & \gate{\mathrm{T^{\dagger}}}    & \qw      & \qw                         & \ctrl{2} & \gate{\mathrm{T}}     & \targ    & \qw & \\
%%                              &            &     &          &  & & \push{y}                  &          & \push{x_a \oplus y}            &          & \push{x_a \oplus x_b \oplus y} &          & \push{x_b \oplus y}         &          &                       &          &     & \\
%%           \lstick{\ket{y}}   & \ctrl{-1}  & \qw &          &  & & \gate{\mathrm{T}}         & \targ    & \gate{\mathrm{T^{\dagger}}}    & \targ    & \gate{\mathrm{T}}              & \targ    & \gate{\mathrm{T^{\dagger}}} & \targ    & \qw                   & \qw      & \qw & \\
%% }

\begin{quantikz}[row sep=0.2em, column sep=0.5em]
\setwiretype{n}        &           &     &                 &          &&                 &     &           & \push{x_a}            &           &          &          & \push{x_a \oplus x_c}  &           &          &             &                                  &             &           &           &       \\
\lstick{$\ket{x_a}$}   & \ctrl{4}  & \qw & \setwiretype{n} &          && \setwiretype{q} & \qw & \qw       & \gate{T}              & \targ{}   & \qw      & \ctrl{2} & \gate{T^\dagger}       & \targ{}   & \ctrl{4} & \qw         & \qw                              & \targ{}     & \qw       & \ctrl{2} & \qw \\
\setwiretype{n}        &           &     &                 &          &&                 &     &           & \push{x_b \oplus x_c} &           &          &          & \push{x_a \oplus x_b } &           &          &             & \push{x_a \oplus x_b \oplus x_c} &             &           &           &       \\
\lstick{$\ket{x_b}$}   & \ctrl{2}  & \qw & \setwiretype{n} &\push{=}  && \setwiretype{q} & \qw & \targ{}   & \gate{T^\dagger}      & \qw       & \ctrl{2} & \targ{}  & \gate{T^\dagger}       & \ctrl{-2} & \qw      & \targ{}     & \gate{T}                         & \ctrl{-2}   & \targ{}   & \targ{}   & \qw \\
\setwiretype{n}        &           &     &                 &          &&                 &     &           & \push{x_c}            &           &          &          & \push{x_b}             &           &          &             &                                  &             &           &           &       \\
\lstick{$\ket{x_c}$}   & \ctrl{-2} & \qw & \setwiretype{n} &          && \setwiretype{q} & \qw & \ctrl{-2} & \gate{T}              & \ctrl{-4} & \targ{}  & \qw      & \gate{T}               & \qw       & \targ{}  & \ctrl{-2}   & \qw                              & \qw         & \ctrl{-2} & \qw       & \qw
\end{quantikz}

    }
    \subcaption{CCZ in CNOT+T}
    \label{fig-cnot-t-ccz}
  \end{minipage}
  \begin{minipage}{\linewidth}
    \centering
    \scalebox{0.7} {
      \input{img/fig-toff-mark-matroid-alt}
    }
    \subcaption{Alternate CCZ in CNOT+T}
    \label{fig-alt-cnot-t-ccz}
  \end{minipage}
  \caption{Two separate CCZ gate in CNOT+T basis}
\end{figure}

From~\figref{fig-cnot-t-ccz}, we can read the actions of the T-gates into sum-of-paths expression as follows: the
first T gate in the first column adds the term $1/4 \cdot x_0$ and adding the second T$^{\dagger}$-gate in the
first column adds the term $-1/4 \cdot (x_1 \oplus x_2)$ (notice how the coefficient has a minus sign because of
the dagger) making the sum-of-paths $1/4 \cdot x_0 -1/4 \cdot (x_1 \oplus x_2)$. We complete these for the
rest of the T gates and we get the following sum-of-paths expression:

\begin{equation}
  \label{eq-toff-bool}
  \begin{aligned}
    &x_a \cdot x_b \cdot x_c = \frac{1}{4}x_a + \frac{1}{4}x_b + \frac{1}{4}x_c - \frac{1}{4}(x_a \oplus x_b) \\
    &\qquad -\frac{1}{4}(x_a \oplus x_c) - \frac{1}{4}(x_b \oplus x_c) + \frac{1}{4}(x_a \oplus x_b \oplus x_c)
  \end{aligned}
\end{equation}

That is, the phase is $\pi$ when $x_a$, $x_b$, and $x_c$ are true. From here, we can verify
that~\figref{fig-cnot-t-ccz} implements the CCZ gate.

Observe that because the CCZ gate is its own inverse, adding~\eqnref{eq-toff-bool} to itself gives us $2\pi = 0$.
In the same way, complementing all of the signs in the sum-of-paths also gives us an expression that can be added
to~\eqnref{eq-toff-bool}, depicted in~\eqnref{eq-toff-bool-alt}. This gives us an alternate way of implementing
the CCZ depicted in~\figref{fig-alt-cnot-t-ccz}.

\begin{equation}
  \label{eq-toff-bool-alt}
  \begin{aligned}
    &x_a \cdot x_b \cdot x_c = - \frac{1}{4}x_a - \frac{1}{4}x_b - \frac{1}{4}x_c + \frac{1}{4}(x_a \oplus x_b) \\
    &\qquad + \frac{1}{4}(x_a \oplus x_c) + \frac{1}{4}(x_b \oplus x_c) - \frac{1}{4}(x_a \oplus x_b \oplus x_c)
  \end{aligned}
\end{equation}

Suppose in our example, the Hadamard gates conjugate the CCZ gate on the $x_c$ qubit. Observe that three terms:
$x_a$, $x_b$, and $x_a \oplus x_b$ have no dependence on $x_c$. Removing them means the Hadamard gates transform
the resulting gate into a relative phase Toffoli gate (RTOF), which implements the Toffoli gate up to relative
phase, where the relative phase is equal to the complement of the terms we eliminated~\cite{bib-amy-phase-state}.
\figref{fig-rtof} shows one implementation of such an RTOF, which can be implemented with three fewer T-gates
than the Toffoli gate. In general, we can use this principle to create relative phase versions of Boolean logic
with lower T-count than the normal versions, as we will seein~\secref{Pro}.

\begin{figure}
  \scalebox{0.8} {
    %% \begin{quantikz}[row sep=0.25cm, column sep=0.35cm]
%% \lstick{} & \ctrl{2}  & \lstick{} &\ghost{\qw}         & \ghost{\qw} & \qw & \qw                & \targ{}   & \qw      & \ctrl{1} & \gate{T}           & \ctrl{1} & \qw      & \targ{}    & \qw \\
%% \lstick{} & \ctrl{1}  & \lstick{} &\push{=} & \ghost{\qw} & \qw & \qw                & \qw       & \ctrl{1} & \targ{}  & \gate{T^{\dagger}} & \targ{}  & \ctrl{1} & \qw        & \qw \\
%% \lstick{} & \gate{iZ} & \lstick{} &\push{ } & \ghost{\qw} & \qw & \gate{T^{\dagger}} & \ctrl{-2} & \targ{}  & \qw      & \gate{T}           & \qw      & \targ{}  & \ctrl{-2}  & \qw 
%% \end{quantikz}
\begin{quantikz}[row sep=0.25cm, column sep=0.35cm]
\lstick{$\ket{x_a}$} & \ctrl{1}                & \setwiretype{n} &          && \setwiretype{q} & \qw &          & \qw & \qw                & \targ{}   & \qw      & \ctrl{1} & \gate{T}           & \ctrl{1} & \qw      & \targ{}    &          & \qw \\
\lstick{$\ket{x_b}$} & \ctrl{1}                & \setwiretype{n} & \push{=} && \setwiretype{q} & \qw &          & \qw & \qw                & \qw       & \ctrl{1} & \targ{}  & \gate{T^{\dagger}} & \targ{}  & \ctrl{1} & \qw        &          & \qw \\
\lstick{$\ket{x_c}$} & \targ[style={double}]{} & \setwiretype{n} &          && \setwiretype{q} & \qw & \gate{H} & \qw & \gate{T^{\dagger}} & \ctrl{-2} & \targ{}  & \qw      & \gate{T}           & \qw      & \targ{}  & \ctrl{-2}  & \gate{H} & \qw \\
\end{quantikz}


  }
  \caption{A relative phase Toffoli gate}
  \label{fig-rtof}
\end{figure}

Observe that because there is a second, complementary implementation of the Toffoli gate, which uses the
complementary sum-of-paths in~\eqnref{eq-toff-bool-alt} we can also create an RTOF from this sum-of-paths
with a complementary relative phase to~\figref{fig-rtof}. We can use these two in tandem later to reduce the
T-count in decompositions such as in~\secref{sub-mct}.


\subsection{Multiple Control Toffoli Gate Decompositions}
\label{sub-mct}

\begin{figure}[t]
  \centering
  \begin{minipage}{\linewidth}
    \centering
    \scalebox{1.0} {
      \begin{quantikz}[row sep=0.5em, column sep=0.5em]
  \setwiretype{n}    &     &          &     &                 &                & &                 & \push{g_0} & \push{g_1} & \push{g_2} &     &                    \\
  \lstick{\ket{x_0}} & \qw & \ctrl{1} & \qw & \setwiretype{n} & \setwiretype{n}& & \setwiretype{q} & \ctrl{1}   & \qw        & \ctrl{1}   & \qw & \rstick{\ket{x_0}} \\
  \lstick{\ket{x_1}} & \qw & \ctrl{1} & \qw & \setwiretype{n} & \setwiretype{n}& & \setwiretype{q} & \ctrl{2}   & \qw        & \ctrl{2}   & \qw & \rstick{\ket{x_1}} \\
  \lstick{\ket{x_2}} & \qw & \ctrl{2} & \qw & \setwiretype{n} & \push{=}       & & \setwiretype{q} & \qw        & \ctrl{1}   & \qw        & \qw & \rstick{\ket{x_2}} \\
  \lstick{\ket{0}}   & \qw & \qw      & \qw & \setwiretype{n} & \setwiretype{n}& & \setwiretype{q} & \targ{}    & \ctrl{1}   & \targ{}    &     & \rstick{\ket{0}} \\  
  \lstick{\ket{y}}   & \qw & \targ{}  & \qw & \setwiretype{n} & \setwiretype{n}& & \setwiretype{q} & \qw        & \targ{}    & \qw        & \qw & \rstick{\ket{x_0 \cdot x_1 \cdot x_2 \oplus y}}  \\
\end{quantikz}                        


    }
    \subcaption{Clean Ancilla Decomposition of MCT}
    \label{fig-clean-mct}
  \end{minipage}
  \begin{minipage}{\linewidth}
    \centering
    \scalebox{1.0} {
      \begin{quantikz}[row sep=0.5em, column sep=0.5em]
  \setwiretype{n}    &     &          &     &                 &                 &&                 & \push{g_0} & \push{g_1} & \push{g_2} & \push{g_3} &  \\  
  \lstick{\ket{x_0}} & \qw & \ctrl{1} & \qw & \setwiretype{n} & \setwiretype{n} && \setwiretype{q} & \qw        & \ctrl{1}   & \qw        & \ctrl{1}   & \qw & \rstick{\ket{x_0}} \\
  \lstick{\ket{x_1}} & \qw & \ctrl{1} & \qw & \setwiretype{n} & \setwiretype{n} && \setwiretype{q} & \qw        & \ctrl{2}   & \qw        & \ctrl{2}   & \qw & \rstick{\ket{x_1}} \\
  \lstick{\ket{x_2}} & \qw & \ctrl{2} & \qw & \setwiretype{n} & \push{=}        && \setwiretype{q} & \ctrl{1}   & \qw        & \ctrl{1}   & \qw        & \qw & \rstick{\ket{x_2}} \\
  \lstick{a}         & \qw & \qw      & \qw & \setwiretype{n} & \setwiretype{n} && \setwiretype{q} & \ctrl{1}   & \targ{}    & \ctrl{1}   & \targ{}    &  & \rstick{a} \\  
  \lstick{\ket{y}}   & \qw & \targ{}  & \qw & \setwiretype{n} & \setwiretype{n} && \setwiretype{q} & \targ{}    & \qw        & \targ{}    & \qw        & \qw & \rstick{\ket{x_0 \cdot x_1 \cdot x_2 \oplus y}}  \\
\end{quantikz}                        


    }
    \subcaption{Dirty Ancilla Decomposition of MCT}
    \label{fig-dirty-mct}
  \end{minipage}
  \caption{MCT Decompositions}
  \label{fig-mct-decomp}
\end{figure}


\begin{figure}[t]
  \scalebox{0.8} {
    \begin{quantikz}[row sep=0.3cm, column sep=0.55cm]
\lstick{$\ket{x_5}$} & \ctrl{1} & \qw      & \qw      & \qw      & \qw      & \qw      & \qw      & \ctrl{1}    \\
\lstick{$\ket{x_4}$} & \ctrl{4} & \qw      & \qw      & \qw      & \qw      & \qw      & \qw      & \ctrl{4}    \\
\lstick{$\ket{x_3}$} & \qw      & \ctrl{3} & \qw      & \qw      & \ctrl{2} & \qw      & \ctrl{3} & \qw          \\
\lstick{$\ket{x_2}$} & \qw      & \qw      & \ctrl{3} & \qw      & \ctrl{4} & \ctrl{3} & \qw      & \qw          \\
\lstick{$\ket{x_1}$} & \qw      & \qw      & \qw      & \ctrl{3} & \qw      & \qw      & \qw      & \qw          \\
\lstick{$a_0$}       & \targ{}  & \ctrl{1} & \qw      & \qw      & \qw      & \qw      & \ctrl{1} & \targ{}           \\
\lstick{$a_1$}       & \qw      & \targ{}  & \targ{}  & \ctrl{1} & \qw      & \targ{}  & \targ{}  & \qw               \\
\lstick{}            & \qw      & \qw      & \qw      & \targ{}  & \targ{}  & \qw      & \qw      & \qw           \\
\end{quantikz}

  }
  \caption{Quantum Circuit for~\figref{fig-xag}}
  \label{fig-xag-qc} 
\end{figure}

Multiple Control Toffoli (MCT) gates are generalized Toffoli gates with an arbitrary number
of control bits. In general decomposing them takes exponential cost~\cite{bib-maslov-rtof}
without any ancilla, while using ancilla allows for a linear
cost~\cite{bib-barenco-elementary}. There are two types of ancilla available for use:
\emph{clean ancilla}, which are initialized to 0, and \emph{dirty ancilla}, which have no
assumptions on their values. While it is widely known that clean ancilla decompositions
can be implemented with fewer Toffoli gates than dirty ancilla
decompositions~\cite{bib-mike-and-ike} in the reversible case , several results since have
demonstrated that using relative phase constructions can bring down the
cost~\cite{bib-maslov-rtof, bib-amy-phase-state, bib-abdessaied-mct}. In fact, it has
recently been demonstrated that linear cost can be achieved using a single
ancilla~\cite{bib-zindorf-mct}. However, all of these methods still show a linear
separation between the cost of decomposing an MCT using clean and dirty ancilla.

We use the dirty ancilla MCT deconstructions as a base for our dirty ancilla XAG
construction. We then demonstrate that in the quantum case, implementing such a dirty
ancilla XAG construction actually has the same cost as the clean ancilla deconstruction.

\subsection{Logic Networks and XOR-And-Inverter Graphs}
\label{sub-xag}
\begin{figure}
  \centering
  \scalebox{1.0} {
    \begin{tikzpicture}[
    gate/.style={circle, draw, inner sep=2pt},
    every node/.style={font=\footnotesize},
    edgelabel/.style={font=\footnotesize, fill=white, inner sep=1pt},
    >=Stealth,
    xstep=1cm, ystep=1cm
]

% ------------------------------------------------------------
% Inputs (bottom layer)
% ------------------------------------------------------------
\node (x1) at (0,0) {$x_1$};
\node (x2) at (1,0) {$x_2$};
\node (x3) at (2,0) {$x_3$};
\node (x4) at (3,0) {$x_4$};
\node (x5) at (4,0) {$x_5$};

% ------------------------------------------------------------
% Gates
% ------------------------------------------------------------
\node[gate] (xorTop) at (0,2) {$+$};
\node[gate] (and1)   at (0,1) {$\wedge$};
\node[gate] (and2)   at (1,2) {$\wedge$};

\node[gate] (xor2) at (1,1) {$+$};
\node[gate] (and3) at (2,1) {$\wedge$};
\node[gate] (and4) at (3,1) {$\wedge$};

% Output
\node (F) at (0,3) {$F$};

% ------------------------------------------------------------
% Edges + fully angled labels
% ------------------------------------------------------------

% xorTop �� F
\draw[->] (xorTop) -- (F);

% and1 �� xorTop (angle upward-left)
\draw (and1) -- (xorTop);
%    node[edgelabel, sloped, above, pos=0.45] {$x_1(\,x_2 \oplus x_3 x_4 x_5\,)$};

% x1 �� and1
\draw (x1) -- (and1);

% and2 �� xorTop (sloped)
\draw (and2) -- (xorTop);
%    node[edgelabel, sloped, above, pos=0.55] {$x_2 x_3$};

% x2,x3 �� and2
\draw[bend left=45] (x2) to (and2);
\draw (x3) -- (and2);
%\draw[bend right=45] (x3) to (and2);

% xor2 �� and1 (label angled with edge)
\draw (xor2) -- (and1);
%    node[edgelabel, sloped, above, pos=0.55] {$x_2 \oplus x_3 x_4 x_5$};

\draw (x2) -- (xor2);

% and3 �� xor2 (sloped)
\draw (and3) -- (xor2);
%    node[edgelabel, sloped, above, pos=0.55] {$x_3 x_4 x_5$};

% and4 �� and3 (sloped)
\draw (and4) -- (and3);
%    node[edgelabel, sloped, above, pos=0.55] {$x_4 x_5$};

% x3 �� and3
\draw (x3) -- (and3);

% x4,x5 �� and4
\draw (x4) -- (and4);
\draw (x5) -- (and4);

\end{tikzpicture}

%% \documentclass[tikz, border=5pt]{standalone}

%% \usepackage{tikz}
%% \usetikzlibrary{arrows.meta, calc, positioning}

%% \begin{document}

%% \begin{tikzpicture}[
%%     gate/.style={circle, draw, inner sep=2pt},
%%     every node/.style={font=\scriptsize},
%%     edgelabel/.style={font=\scriptsize, fill=white, inner sep=1pt},
%%     >=Stealth
%% ]

%% % Inputs
%% \node (x1) at (0,0) {$x_1$};
%% \node (x2) at (1.5,0) {$x_2$};
%% \node (x3) at (3.0,0) {$x_3$};
%% \node (x4) at (4.5,0) {$x_4$};
%% \node (x5) at (6.0,0) {$x_4$};

%% % Output arrow to F
%% \node (F) at (0,7.0) {$F$};

%% % Top XOR
%% \node[gate] (xorTop) at (0,6.0) {$+$};
%% \draw[->] (xorTop) -- (F);

%% % AND gate
%% \node[gate] (and1) at (0,4.5) {$\wedge$};
%% \draw (and1) -- (xorTop)
%%     node[edgelabel, pos=0.5, xshift=-15pt]{$x_1 (x_2 \oplus x_3 x_4 x_5)$};
%% \draw (x1) -- (and1);

%% % x2x3
%% \node[gate] (and2) at (1.5,6.0) {$\wedge$};
%% \draw (and2) -- (xorTop)
%%     node[edgelabel, pos=0.5, xshift=-15pt]{$x_2 x_3$};
%% \draw (x2) -- (and2);
%% \draw (x3) -- (and2);

%% % x2 ^ x3x4x5
%% \node[gate] (xor2) at (1.5,4.5) {$+$};
%% \draw (xor2) -- (and1)
%%     node[edgelabel, pos=0.5, xshift=-15pt]{$x_2 \oplus x_3 x_4 x_5$};
%% \draw (x2) -- (and2);

%% % x2x3
%% \node[gate] (and3) at (3.0, 3.0) {$\wedge$};
%% \draw (and3) -- (xor2)
%%     node[edgelabel, pos=0.5, xshift=-15pt]{$x_2 x_3$};
%% \draw (x3) -- (and2);

%% % x2x3
%% \node[gate] (and4) at (4.5, 1.5) {$\wedge$};
%% \draw (and4) -- (and3)
%%     node[edgelabel, pos=0.5, xshift=-15pt]{$x_2 x_3$};
%% \draw (x4) -- (and4);
%% \draw (x5) -- (and4);    

%% \end{tikzpicture}

%% \end{document}

  }
  \caption{An XOR-AND-Inverter Graph for $x_1(x_2 \oplus x_3x_4)\oplus x_2x_3$}
  \label{fig-xag}
\end{figure}

A {\it Boolean logic network} is a simple Directed Acyclic Graph (DAG). It has nodes,
edges, primary outputs, and primary inputs. Each Boolean logic network represents a
Boolean function, and a \emph{node} represents a subfunction used to decompose that Boolean
function. It has zero or more {\it fanins}, or incoming edges, as well as zero or more
{\it fanouts} or outgoing edges. The {\it primary inputs (PIs)} are nodes without fanins
in the current network. The {\it primary outputs} are a subset of nodes in the network
without fanouts.

XOR-AND-Inverter graphs (XAG) are examples of such logic networks, where each node
represents either the AND operation $\land$ or the XOR operation $\oplus$, while
bubbles on either the input fanins or the output fanins represent negations of those
values.~\figref{fig-xag} shows an example of such a XAG.

\subsection{Synthesizing Quantum Circuits from XAG}
\label{sub-synth-xag}
Meuli et al~\cite{bib-meuli-ros} demonstrated that quantum circuits can be generated from
XAG. AND nodes are implemented by routing the fanins to the input of a Toffoli gate that
acts on a clean ancilla. XOR nodes are implemented by placing the targets of the
subcircuits implementing the fanin nodes onto a clean ancilla in series; any fanins that
are primary inputs are implemented as a CNOT gate with its control on the primary input
and the target on the same clean ancilla that the other fanin acts on. Any actions on
clean ancilla needs to be uncomputed.

Because the AND nodes are implemented using Toffoli gates, the number of AND nodes in
the circuit corresponds directly to the T-count of the circuit. Methods to
optimize XAG for these metrics thus present an opportunity to optimize for T-count in
these methods.

Converting an XAG to a quantum circuit is straightforward exercise. Any AND node is 
turned into a Toffoli gate, with its fanins routed to the Toffoli gate's inputs, and
its target set to a clean ancilla, which describes the AND node's output. XOR operations
are implemented with the incoming subcircuits acting on the same clean ancilla.

While there are many ways of traversing this XAG to optimize for ancilla qubits, we
focus on the \emph{Bennett method}~\cite{bib-time-space-bennett}, which traverses the
XAG depth first from the inputs in order to minimize the number of times a value has to
be uncomputed. The result of generating such an XAG is depicted in~\figref{fig-xag-qc}.

We can apply the same replacement with relative phase Toffoli as we did in~\secref{sub-mct}
in order to achieve a lower T-count. 




% \input{Mot}

\section{Proposed}
\label{Pro}

\subsection{Dirty Ancilla Node Gadgets}
\label{sub-synth}
Here we take the dirty ancilla construction in~\ref{sub-mct}, and create a gadget to
realize an XAG AND node using dirty ancilla. First, we separate the construction into
two parts, as in~\figref{fig-dirt-mct-org}: we label these \emph{compute}, and \emph{uncompute}.
We observe that the Toffoli gate in compute can be replaced with a single target controlled
gate for an arbitrary function F; that is, it's a gate that flips its target bit when a
function F is true. We can implement this function F as a quantum
subcircuit~\cite{bib-amy-phase-state}, and so we can realize a recursive construction, depicted
in~\figref{fig-compute}. The Toffoli gate in the uncompute gadget can thus be replaced with the
gate F's adjoint. For a relative phase gate, this means implementing F's components' adjoints in
reverse, to realize a gate that implements the same function, but with the opposite relative phase
~\cite{bib-maslov-rtof}.

\begin{figure}[t]
  \centering
  \scalebox{1.0} {
    \begin{quantikz}[row sep=0.5em, column sep=0.5em]
  \setwiretype{n}    &     &                                                                                                                                                                 & \push{Compute} &          &  & \push{Uncompute}                                                                                                                                                     &  &  \\
  \setwiretype{n}    &     &                                                                                                                                                                 &                &          &  &                                                                                                                                                                      &  &  \\  
  \lstick{\ket{x_0}} & \qw & \qw \gategroup[5,steps=3,style={dashed,rounded corners,fill=blue!20, inner xsep=2pt},background,label style={label position=above,anchor=north,yshift=-0.2cm}]{}& \ctrl{1}       & \qw      &  & \ctrl{1} \gategroup[5,steps=2,style={dashed,rounded corners,fill=red!20, inner xsep=2pt},background,label style={label position=above,anchor=north,yshift=-0.2cm}]{} &  & \rstick{\ket{x_0}} \\
  \lstick{\ket{x_1}} & \qw & \qw                                                                                                                                                             & \ctrl{2}       & \qw      &  & \ctrl{2}                                                                                                                                                             &  & \rstick{\ket{x_1}} \\
  \lstick{\ket{x_2}} & \qw & \ctrl{1}                                                                                                                                                        & \qw            & \ctrl{1} &  & \qw                                                                                                                                                                  &  & \rstick{\ket{x_2}} \\
  \lstick{a}         & \qw & \ctrl{1}                                                                                                                                                        & \targ{}        & \ctrl{1} &  & \targ{}                                                                                                                                                              &  & \rstick{a} \\  
  \lstick{\ket{y}}   & \qw & \targ{}                                                                                                                                                         & \qw            & \targ{}  &  & \qw                                                                                                                                                                  &  & \rstick{\ket{x_0 \cdot x_1 \oplus y}}  
\end{quantikz}

%% \begin{quantikz}[row sep=0.5em, column sep=0.5em]
%%   \lstick{\ket{x_0}} & \qw & \qw      & \ctrl{1} & \qw      & \ctrl{1} & \qw & \rstick{\ket{x_0}} \\
%%   \lstick{\ket{x_1}} & \qw & \qw      & \ctrl{2} & \qw      & \ctrl{2} & \qw & \rstick{\ket{x_1}} \\
%%   \lstick{\ket{x_2}} & \qw & \ctrl{1}\gategroup[5,steps=3,style={dashed,rounded corners},label style={above}]{Compute} & \qw      & \ctrl{1} & \qw      & \qw & \rstick{\ket{x_2}} \\
%%   \lstick{a}         & \qw & \ctrl{1} & \targ{}  & \ctrl{1} & \targ{}  &     & \rstick{a} \\  
%%   \lstick{\ket{y}}   & \qw & \targ{}  & \qw      & \targ{}\gategroup[5,steps=1,style={dashed,rounded corners},label style={above}]{Uncompute} & \qw      & \qw & \rstick{\ket{x_0 \cdot x_1 \oplus y}}  \\
%% \end{quantikz}
%% \begin{quantikz}[row sep=0.5em, column sep=0.5em]
%%   \lstick{\ket{x_0}} & \qw & \qw      & \ctrl{1} & \qw      & \ctrl{1} & \qw & \rstick{\ket{x_0}} \\
%%   \lstick{\ket{x_1}} & \qw & \qw      & \ctrl{2} & \qw      & \ctrl{2} & \qw & \rstick{\ket{x_1}} \\
%%   \lstick{\ket{x_2}} & \qw & \ctrl{1} & \qw      & \ctrl{1} & \qw      & \qw & \rstick{\ket{x_2}} \\
%%   \lstick{a}         & \qw & \ctrl{1} & \targ{}  & \ctrl{1} & \targ{}  &  & \rstick{a} \\  
%%   \lstick{\ket{y}}   & \qw & \targ{}  & \qw      & \targ{}  & \qw      & \qw & \rstick{\ket{x_0 \cdot x_1 \oplus y}}  \\
%% \end{quantikz}                        


  }
  \caption{Dirty MCT decomposition base}
  \label{fig-dirt-mct-org}
\end{figure}

\begin{figure}[t]
  \centering
  \scalebox{1.0} {
    \begin{quantikz}[row sep=0.3cm, column sep=0.3cm]
\qw              & \qw             & \qw        & \gate[2]{F}         & \qw               & \qw      & \qw    \\
\qw              & \qw             & \qw        & \ghost{F}\ctrl{2}   & \qw               & \qw      & \qw    \\
\lstick{$G$}     & \qw             & \ctrl{1}   & \qw                 & \ctrl{1}          & \qw      & \qw    \\
\lstick{$a_0$}   & \qw             & \gate{Z}   & \targ{}             & \gate{Z^\dagger}  & \qw      & \qw    \\
\lstick{$t$}     & \gate{H}        & \ctrl{-1}  & \qw                 & \ctrl{-1}         & \gate{H} & \qw    \\
\end{quantikz}

%% \begin{quantikz}[row sep=0.3cm, column sep=0.55cm]
%% \qw              & \qw             & \qw        & \qw             & \qw              & \gate[2]{F_i}       & \qw                    & \qw       & \qw             & \qw      & \qw        & \qw              & \gate[2]{F_i}        & \qw                     & \qw        \\
%% \qw              & \qw             & \qw        & \qw             & \qw              & \ghost{F_i}         & \qw                    & \qw       & \qw             & \qw      & \qw        & \qw              & \ghost{F_i}          & \qw                     & \qw        \\
%% \qw              & \qw             & \qw        & \qw             & \ctrl{2}         & \qw                 & \ctrl{2}               & \qw       & \qw             & \qw      & \qw        & \ctrl{2}         & \qw                 & \ctrl{2}                & \qw         \\
%% \qw              & \qw             & \octrl{2}  & \qw             & \qw              & \qw                 & \qw                    & \qw       & \qw             & \qw      & \qw        & \qw              & \qw                 & \qw                     & \qw         \\
%% \lstick{$a_1$}   & \qw             & \qw        & \qw             & \gate{i\omega Z} & \targ{}             & \gate{i\omega Z^\dagger} & \qw       & \octrl{1}       & \qw      & \qw        & \gate{i\omega Z} & \targ{}             & \gate{i\omega Z^\dagger}  & \qw          \\
%% \lstick{$a_2$}   & \qw             & \gate{iZ}  & \gate{H}        & \ctrl{-1}        & \qw                 & \ctrl{-1}              & \gate{H}  & \gate{iZ^\dagger} & \qw      & \gate{H}   & \ctrl{-1}        & \qw                 & \ctrl{-1}               & \gate{H}     \\
%% \lstick{}        & \gate{H}        & \ctrl{-1}  & \qw             & \qw              & \qw                 & \qw                    & \qw       & \ctrl{-1}       & \gate{H} & \qw        & \qw              & \qw                 & \qw                     & \qw          \\
%%\end{quantikz}


  }
  \caption{Compute gadget}
  \label{fig-compute}
\end{figure}

We utilize the fact that the CCZ has two complementary implementations to our advantage by replacing the left CCZ
gate with~\figref{fig-cnot-t-ccz} and the right with \figref{fig-alt-cnot-t-ccz}. In so doing we can combine their
sum-of-paths to get the below sum-of-paths equation:

\begin{equation}
  \label{eq-toff-tot}
  \begin{aligned}
    &\frac{1}{4}z + \frac{1}{4}a + \frac{1}{4}f - \frac{1}{4}(z \oplus a) \\
    &\qquad -\frac{1}{4}(z \oplus f) - \frac{1}{4}(a \oplus f) + \frac{1}{4}(z \oplus a \oplus f)\\
    &- \frac{1}{4}z - \frac{1}{4}(xy \oplus a) - \frac{1}{4} f + \frac{1}{4}(z \oplus xy \oplus a) \\
    &\qquad + \frac{1}{4}(z \oplus f) + \frac{1}{4}(xy \oplus a \oplus f) - \frac{1}{4}(z \oplus xy \oplus a \oplus f)
  \end{aligned}
\end{equation}

This causes some cancellations, and we are left with:

\begin{equation}
  \label{eq-toff-simp}
  \begin{aligned}
    &\frac{1}{4}a - \frac{1}{4}(z \oplus a) - \frac{1}{4}(xy \oplus a) + \frac{1}{4}(z \oplus xy \oplus a)  \\
    &+ \frac{1}{4}(z \oplus a \oplus f) - \frac{1}{4}(a \oplus f) + \frac{1}{4}(xy \oplus a \oplus f) \\
    &- \frac{1}{4}(z \oplus xy \oplus a \oplus f)
  \end{aligned}
\end{equation}

Observe that the first line in\eqnref{eq-toff-simp} contains only terms that have no dependence on $f$.
That means we can bring them out from the summation to create a relative phase version of the compute
gadget, which we depict in~\figref{fig-dirty-gadget-rp}. This means that we can implement the compute
gadget with the same T-count as a relative phase Toffoli gate. We can similarly formulate an uncompute
gadget that is this gadget's conjugate, in order to uncompute the value for the same cost. This implies
that any intermediate computation using dirty ancilla can be performed with the same cost as a relative
phase Toffoli gate. 

The output node will have to implement all 8 terms in~\eqnref{eq-toff-simp}. This means that the T-count of
implementing the output node gadget is 8, so that the difference in T-count between the clean ancilla and
the dirty ancilla decompositions is dependent on the number of these output nodes, which is the number of
output bits of the Boolean function $m$. This means that for an MCT gate, the difference between clean
ancilla and dirty ancilla is 1 T-gate. We can implement these as two RTOF gates.

\begin{table*}[t]
\centering
\scalebox{0.9}{
  \begin{tabular}{p{0.05\textwidth}|p{0.10\textwidth}|p{0.10\textwidth}|p{0.75\textwidth}}
    \# & Node & Conditions & Gadget \\\hline
       {1}
       &
       \scalebox{1.0} {
         \input{xag_and}
       }
       &
       One of A or B are primary inputs and the other is a fanin from another node.
       All fanouts F are intermediate fanouts
       &
       \scalebox{0.9} {
         \input{fig-prop-and}
       }\\\hline
       {2}
       &
       \scalebox{1.0} {
         \input{xag_and}
       }
       &
       Both A and B are fanins from other nodes
       &
       \scalebox{0.9} {
         \begin{quantikz}[row sep=0.2em, column sep=0.5em]
\qw                & \gate[2]{B}              & \gate[2]{B}              & \qw             &                    & \setwiretype{n} & \setwiretype{n} & \setwiretype{n} & \setwiretype{q} & \qw           & \gate[2]{B}       & \qw          & \qw           & \qw                & \qw           & \gate[2]{A}            & \qw           & \qw                & \qw           & \qw            & \qw           & \qw                     \\
\qw                & \ghost{B}\ctrl{2}        & \ghost{B}\ctrl{2}        & \qw             &                    & \setwiretype{n} & \setwiretype{n} & \setwiretype{n} & \setwiretype{q} & \qw           & \ghost{B}\ctrl{2} & \qw          & \qw           & \qw                & \qw           & \ghost{A}\ctrl{2}      & \qw           & \qw                & \qw           & \qw            & \qw           & \qw                     \\
\setwiretype{n}    & \setwiretype{n}          & \setwiretype{n}          & \setwiretype{n} &                    & \push{=}        &                 &                 & \setwiretype{n} & \push{\vdots} & \setwiretype{n}   & \push{\vdots}& \push{\vdots} & \push{\vdots}      & \push{\vdots} & \push{\vdots}          & \push{\vdots} & \push{\vdots}      & \push{\vdots} & \push{\vdots}  & \push{\vdots} & \push{\vdots}               \\
\lstick{$\ket{0}$} & \targ{}                  & \gate[1]{\wedge}\ctrl{2} & \qw             & \push{B}           & \setwiretype{n} & \setwiretype{n} & \setwiretype{n} & \setwiretype{q} & \qw           & \targ{}           & \qw          & \targ{}       & \gate{T}           & \ctrl{1}      & \targ[style={double}]{}& \ctrl{1}      & \gate{T^{\dagger}} & \targ{ }      & \qw            & \qw           & \qw   \\
\lstick{$a$}       &                          & \targX{}                 & \qw             & \push{A \oplus a}  & \setwiretype{n} & \setwiretype{n} & \setwiretype{n} & \setwiretype{q} & \qw           &                   & \targ{}      & \ctrl{-1}     & \gate{T^{\dagger}} & \targ{}       & \qw                    & \targ{}       & \gate{T}           & \ctrl{-1}     & \targ{}        & \qw           & \qw   \\
\lstick{$f$}       &                          & \targ[style={double}]{}  & \qw             & \push{AB \oplus f} & \setwiretype{n} & \setwiretype{n} & \setwiretype{n} & \setwiretype{q} & \gate{H}      &                   & \ctrl{-1}    & \qw           & \qw                & \qw           & \qw                    & \qw           & \qw                & \qw           & \ctrl{-1}      & \gate{H}      & \qw         
\end{quantikz}

%% \begin{quantikz}[row sep=0.3cm, column sep=0.55cm]
%% \qw              & \qw             & \qw        & \qw             & \qw              & \gate[2]{F_i}       & \qw                    & \qw       & \qw             & \qw      & \qw        & \qw              & \gate[2]{F_i}        & \qw                     & \qw        \\
%% \qw              & \qw             & \qw        & \qw             & \qw              & \ghost{F_i}         & \qw                    & \qw       & \qw             & \qw      & \qw        & \qw              & \ghost{F_i}          & \qw                     & \qw        \\
%% \qw              & \qw             & \qw        & \qw             & \ctrl{2}         & \qw                 & \ctrl{2}               & \qw       & \qw             & \qw      & \qw        & \ctrl{2}         & \qw                 & \ctrl{2}                & \qw         \\
%% \qw              & \qw             & \octrl{2}  & \qw             & \qw              & \qw                 & \qw                    & \qw       & \qw             & \qw      & \qw        & \qw              & \qw                 & \qw                     & \qw         \\
%% \lstick{$a_1$}   & \qw             & \qw        & \qw             & \gate{i\omega Z} & \targ{}             & \gate{i\omega Z^\dagger} & \qw       & \octrl{1}       & \qw      & \qw        & \gate{i\omega Z} & \targ{}             & \gate{i\omega Z^\dagger}  & \qw          \\
%% \lstick{$a_2$}   & \qw             & \gate{iZ}  & \gate{H}        & \ctrl{-1}        & \qw                 & \ctrl{-1}              & \gate{H}  & \gate{iZ^\dagger} & \qw      & \gate{H}   & \ctrl{-1}        & \qw                 & \ctrl{-1}               & \gate{H}     \\
%% \lstick{}        & \gate{H}        & \ctrl{-1}  & \qw             & \qw              & \qw                 & \qw                    & \qw       & \ctrl{-1}       & \gate{H} & \qw        & \qw              & \qw                 & \qw                     & \qw          \\
%%\end{quantikz}


       }\\\hline
       {3}
       &
       \scalebox{1.0} {
         \input{xag_and}
       }
       &
       A and B are primary inputs
       &
       \scalebox{1.0} {
         %% \begin{quantikz}[row sep=0.25cm, column sep=0.35cm]
%% \lstick{} & \ctrl{2}  & \lstick{} &\ghost{\qw}         & \ghost{\qw} & \qw & \qw                & \targ{}   & \qw      & \ctrl{1} & \gate{T}           & \ctrl{1} & \qw      & \targ{}    & \qw \\
%% \lstick{} & \ctrl{1}  & \lstick{} &\push{=} & \ghost{\qw} & \qw & \qw                & \qw       & \ctrl{1} & \targ{}  & \gate{T^{\dagger}} & \targ{}  & \ctrl{1} & \qw        & \qw \\
%% \lstick{} & \gate{iZ} & \lstick{} &\push{ } & \ghost{\qw} & \qw & \gate{T^{\dagger}} & \ctrl{-2} & \targ{}  & \qw      & \gate{T}           & \qw      & \targ{}  & \ctrl{-2}  & \qw 
%% \end{quantikz}
\begin{quantikz}[row sep=0.25cm, column sep=0.35cm]
\lstick{$\ket{x_a}$} & \ctrl{1}                & \setwiretype{n} &          && \setwiretype{q} & \qw &          & \qw & \qw                & \targ{}   & \qw      & \ctrl{1} & \gate{T}           & \ctrl{1} & \qw      & \targ{}    &          & \qw \\
\lstick{$\ket{x_b}$} & \ctrl{1}                & \setwiretype{n} & \push{=} && \setwiretype{q} & \qw &          & \qw & \qw                & \qw       & \ctrl{1} & \targ{}  & \gate{T^{\dagger}} & \targ{}  & \ctrl{1} & \qw        &          & \qw \\
\lstick{$\ket{x_c}$} & \targ[style={double}]{} & \setwiretype{n} &          && \setwiretype{q} & \qw & \gate{H} & \qw & \gate{T^{\dagger}} & \ctrl{-2} & \targ{}  & \qw      & \gate{T}           & \qw      & \targ{}  & \ctrl{-2}  & \gate{H} & \qw \\
\end{quantikz}


       }\\\hline
       {4}
       &
       \scalebox{1.0} {
         \input{xag_and}
       }
       &
       F is a primary output or fans out into an XOR node that is a primary output
       &
       \scalebox{0.9} {
         \begin{quantikz}[row sep=0.3cm, column sep=0.55cm]
\qw                & \gate[2]{A}              & \qw             &                     & \setwiretype{n} & \setwiretype{n} & \setwiretype{n} & \setwiretype{q} & \qw                    & \gate[3]{A}       & \qw                    & \qw \\
\qw                & \ghost{A}\ctrl{2}        & \qw             &                     & \setwiretype{n} & \setwiretype{n} & \setwiretype{n} & \setwiretype{q} & \qw                    & \ghost{A}\ctrl{3} & \qw                    & \qw \\
\setwiretype{n}    & \setwiretype{n}          & \setwiretype{n} &                     & \push{=}        &                 &                 & \setwiretype{n} &                        &                   &                        &     \\
\lstick{$B$}       & \gate[1]{\wedge}\ctrl{2} & \qw             & \push{B}            & \setwiretype{n} & \setwiretype{n} & \setwiretype{n} & \setwiretype{q} & \ctrl{1}               & \qw               & \ctrl{1}               & \qw \\
\lstick{$a_i$}     & \targX{}                 & \qw             & \push{A \oplus a_i} & \setwiretype{n} & \setwiretype{n} & \setwiretype{n} & \setwiretype{q} & \ctrl{1}               & \targ{}           & \ctrl{1}               & \qw \\
\lstick{$f$}       & \targ{}                  & \qw             & \push{AB \oplus f}  & \setwiretype{n} & \setwiretype{n} & \setwiretype{n} & \setwiretype{q} & \targ[style={double}]{}& \qw               & \targ[style={double}]{}& \qw  \\
\end{quantikz}

%% \begin{quantikz}[row sep=0.3cm, column sep=0.55cm]
%% \qw              & \qw             & \qw        & \qw             & \qw              & \gate[2]{F_i}       & \qw                    & \qw       & \qw             & \qw      & \qw        & \qw              & \gate[2]{F_i}        & \qw                     & \qw        \\
%% \qw              & \qw             & \qw        & \qw             & \qw              & \ghost{F_i}         & \qw                    & \qw       & \qw             & \qw      & \qw        & \qw              & \ghost{F_i}          & \qw                     & \qw        \\
%% \qw              & \qw             & \qw        & \qw             & \ctrl{2}         & \qw                 & \ctrl{2}               & \qw       & \qw             & \qw      & \qw        & \ctrl{2}         & \qw                 & \ctrl{2}                & \qw         \\
%% \qw              & \qw             & \octrl{2}  & \qw             & \qw              & \qw                 & \qw                    & \qw       & \qw             & \qw      & \qw        & \qw              & \qw                 & \qw                     & \qw         \\
%% \lstick{$a_1$}   & \qw             & \qw        & \qw             & \gate{i\omega Z} & \targ{}             & \gate{i\omega Z^\dagger} & \qw       & \octrl{1}       & \qw      & \qw        & \gate{i\omega Z} & \targ{}             & \gate{i\omega Z^\dagger}  & \qw          \\
%% \lstick{$a_2$}   & \qw             & \gate{iZ}  & \gate{H}        & \ctrl{-1}        & \qw                 & \ctrl{-1}              & \gate{H}  & \gate{iZ^\dagger} & \qw      & \gate{H}   & \ctrl{-1}        & \qw                 & \ctrl{-1}               & \gate{H}     \\
%% \lstick{}        & \gate{H}        & \ctrl{-1}  & \qw             & \qw              & \qw                 & \qw                    & \qw       & \ctrl{-1}       & \gate{H} & \qw        & \qw              & \qw                 & \qw                     & \qw          \\
%%\end{quantikz}


       }\\\hline
       {4}
       &
       \scalebox{1.0} {
         \begin{tikzpicture}[
    gate/.style={circle, draw, inner sep=2pt},
    every node/.style={font=\small}
]

% Inputs
\node (x1) at (0,0) {$A$};
\node (x2) at (1.2,0) {$B$};
\node (x3) at (0.6,1.6) {$F$};

% XOR(x1,x2)
\node[gate] (xor12) at (0.6,0.8) {$+$};
\draw (x1) -- (xor12);
\draw (x2) -- (xor12);
\draw[->] (xor12) -- (x3);

\end{tikzpicture}

       }
       &
       Any XOR node
       &
       \scalebox{1.0} {
         \begin{quantikz}[row sep=0.3cm, column sep=0.55cm]
\qw            & \qw                  & \qw                & \qw   \\
\qw            & \gate[3]{B}          & \gate[3]{A}        & \qw  \\
\qw            & \ghost{B}            & \ghost{A}          & \qw \\
\qw            & \ghost{B}\ctrl{1}    & \ghost{A}\ctrl{1}  & \qw    \\
\qw            & \targ{}              & \targ{}            & \qw  \\
\end{quantikz}
%% \begin{quantikz}[row sep=0.3cm, column sep=0.55cm]
%% \qw            & \qw                            & \gate[3]{A}          & \qw                            & \gate[3]{A}       & \qw   \\
%% \qw            & \gate[3]{B\cdot a_0}         & \ghost{A}            & \gate[3]{B\cdot a_0}         & \ghost{A}         & \qw  \\
%% \qw            & \ghost{B\cdot a_0}           & \ghost{A}\ctrl{1}    & \ghost{B\cdot a_0}           & \ghost{A}\ctrl{1} & \qw \\
%% \lstick{$a_0$} & \ghost{B\cdot a_0}\ctrl{1}   & \targ{}                & \ghost{B\cdot a_0}\ctrl{1}   & \targ{}             & \qw    \\
%% \qw            & \targ{}                        & \qw                    & \targ{}                        & \qw                 & \qw  \\
%% \end{quantikz}


       }
  \end{tabular}
}
\caption{Dirty Ancilla Gadgets}
\label{fig-gadgets}
\end{table*}

%% \begin{table*}[t]
%% \centering
%% \begin{tabular}{p{0.25\textwidth}|p{0.24\textwidth}|p{1.5\textwidth}}
%%   Node & Conditions & Gadget \\\hline
%%   \subfloat[AND-node]{
%%     \makebox[0.25\textwidth][c]{\input{xag_and_node}
%%   }
%%   &
%%   F and G are primary inputs and all fanouts are intermediate fanouts
%%   &
%%   \subfloat[AND-gadget]{
%%     \makebox[\textwidth][c]{\input{fig-prop-and}}
%%   }
%% \end{tabular}
%% \caption{Phase Gadgets}
%% \label{fig-gadgets}
%% \end{table*}


%% The output node will have to implement all 8 terms in~\eqnref{eq-toff-simp}. Observe, however that
%% $\frac{1}{4}a - \frac{1}{4}(z \oplus a) = - \frac{1}{4}z + \frac{1}{2}(a)$. The $\frac{1}{2}$ coefficient
%% can be implemented by an S gate, which is a Clifford gate and has a much more efficient fault-tolerant
%% implementation in most error correcting codes. This means that the T-count of implementing the output
%% node gadget is 7, the exact same as the Toffoli gate which is used as an output node in clean ancilla
%% decomposition.

\begin{figure}
  \scalebox{0.6} {
    \input{fig-prop-and}
  }
  \caption{Relative Phase version of dirty ancilla intermediate compute gadget}
  \label{fig-dirty-gadget-rp}
\end{figure}

\subsection{Proposed Method}

Our proposed method is as follows

\begin{enumerate}
\item Get the next node in the traversal of the XAG, usually from the input
\item Create a subcircuit according to~\tabref{fig-gadgets}
\item Pass on this subcircuit to the next iteration of the algorithm
\item If this node fans out into a primary output, or an XOR node that fans out
  into a primary output take the complement of the passed down subcircuit and set
  it to act on the ancilla that the node uses.
\item If the node is dependent on another subcircuit/subgraph, block until that subcircuit is
  processed.
\item Terminate when all nodes have been processed.
\end{enumerate}

\begin{example}
  \label{ex-dirty-xag}
  We again generate the XAG for~\figref{fig-xag}. Starting from $x_5$ and $x_4$, we
  generate the node they feed into as an RTOF, and pass this RTOF onto the next stage of the iteration.
  We also generate the node taking in $x_2$ and $x_3$, and see that it feeds into an XOR node that
  fans out into a primary output. We generate a Toffoli, and block for the other input.
  We next generate that node's fanout to $x_3$, Because this takes in a fanin from a primary
  input and another node, we plug the RTOF passed down from the previous iteration as $A$
  in gadget\#1, and pass this subcircuit down to the next part of the iteration. We move onto
  the XOR node taking in $x_2$. Since that is an XOR, we generate a subcircuit using a CNOT with
  $x_2$ as an input and replace it in gadget\#4 along with the passed down subcircuit, and pass this
  subcircuit down further. We again repeat the AND procedure for the $x_1$ node the XOR feeds into,
  and feed that output to the next step. Once we have reached the last XOR node, we process both
  the Toffoli feeding into it and the other branch, and generate the full circuit. The full circuit
  is depicted in~\figref{fig-prop-decomp}. We can verify from this diagram that the T-count of this
  is indeed 1 more than the T-count of~\figref{fig-xag-qc} after optimizing it to use RTOF to reduce
  its T-count.
\end{example}

\begin{figure*}[t]
  \centering
  \scalebox{0.8} {
    \begin{quantikz}[row sep=0.2cm, column sep=0.2cm]
\lstick{$x_5$}     & \qw                     & \qw      & \qw             & \qw      &           &                    & \ctrl{1}               &                    &           & \qw       & \qw                      & \qw       & \qw                     & \qw       & \qw             & \qw      &           &                    & \ctrl{1}               &                    &           & \qw       & \qw      & \qw       \\
\lstick{$x_4$}     & \qw                     & \qw      & \qw             & \qw      &           &                    & \ctrl{1}               &                    &           & \qw       & \qw                      & \qw       & \qw                     & \qw       & \qw             & \qw      &           &                    & \ctrl{2}               &                    &           & \qw       & \qw      & \qw       \\
\lstick{$x_3$}     & \qw                     & \qw      & \qw             & \qw      & \targ{}   & \gate{T}           & \targ[style={double}]{}& \gate{T^{\dagger}} & \targ{}   & \qw       & \ctrl{2}                 & \qw       & \qw                     & \ctrl{5}  & \qw             & \qw      & \targ{}   & \gate{T^{\dagger}} & \targ[style={double}]{}& \gate{T}           & \targ{}   & \qw       & \qw      & \qw   \\
\lstick{$x_2$}     & \qw                     & \ctrl{3} & \qw             & \qw      &           &                    & \qw                    &                    &           & \qw       & \qw                      & \qw       & \qw                     & \ctrl{4}  & \qw             & \qw      &           &                    & \qw                    &                    &           & \qw       & \qw      & \ctrl{3}       \\
\lstick{$x_1$}     & \ctrl{1}                &          & \qw             & \qw      &           &                    & \qw                    &                    &           & \qw       & \qw                      & \qw       & \ctrl{1}                & \qw       & \qw             & \qw      &           &                    & \qw                    &                    &           & \qw       & \qw      & \qw      \\
\lstick{$a_1$}     & \qw                     & \qw      & \qw             & \targ{}  & \ctrl{-3} & \gate{T^{\dagger}} & \qw                    & \gate{T}           & \ctrl{-3} & \targ{}   & \gate{i\omega Z^\dagger} & \qw       & \qw                     & \qw       & \qw             & \targ{}  & \ctrl{-3} & \gate{T}           & \qw                    & \gate{T^{\dagger}}  & \ctrl{-3} & \targ{}   & \qw      & \qw       \\
\lstick{$a_2$}     & \ctrl{1}                & \targ{}  & \gate{H}        & \ctrl{-1}&           &                    & \qw                    &                    &           & \ctrl{-1} & \ctrl{-1}                & \gate{H}  & \ctrl{1}                & \qw       & \gate{H}        & \ctrl{-1}&           &                    & \qw                    &                    &           & \ctrl{-1} & \gate{H} & \targ{}   \\
\lstick{$\ket{0}$} & \targ[style={double}]{} & \qw      & \qw             & \qw      &           &                    & \qw                    &                    &           & \qw       & \qw                      & \qw       & \targ[style={double}]{} & \targ{}   & \qw             & \qw      &           &                    & \qw                    &                    &           & \qw       & \qw      & \qw       \\
\end{quantikz}

  }
  \caption{Synthesis of circuit in Example~\ref{ex-dirty-xag}}
  \label{fig-prop-decomp}
\end{figure*}
  
%% \begin{algorithm}[h]
%%   \caption{Proposed Method}
%%   \label{alg-prop}
%%   \begin{algorithmic}[1]

%%     \State $QC \gets$ empty quantum circuit
%%     \State $xag \gets$ XOR-And-Inverter Graph of the function
%%     \State $QC \gets \textsc{parseXAG}(xag, QC)$

%%     \Function{parseXAG}{xag, QC}
%%     \If{$node \gets$ $next\_node(xag)$}
%%        \State $node \gets$ $next\_node(xag)$ \Comment{This is set by a traversal algorithm}
%%     \Else
%%        \State \Return $QC$
%%     \EndIf
    
%%     \Switch{}
%%     \Case{$node$ is XOR}
%%       \State $subQC \gets$ generate \#4
%%     \EndCase
%%     \Case{$node$ is AND and $A$ and $B$ are primary inputs}
%%       \State $subQC \gets$ generate \#2
%%     \EndCase
%%     \Case{$node$ is AND node output is primary output}
%%          \State $subQC \gets$ generate \#3
%%          \State $Uncompute \gets$ QC acting on ancilla bit of subQC
%%          \State $subQC \gets$ Uncompute appended to subQC
%%     \EndCase
%%     \Case{$node$ is AND and one of $A$, $B$ is primary input}
%%          \State $A$ or $B$ $\gets$ QC \Comment{QC is the subcircuit implementing $A$ or $B$}   
%%          \State $subQC \gets$ generate \#1
%%     \EndCase
%%     \EndSwitch
    
%%     \State $next_xags \gets$ subgraphs from $node$'s fanouts
%%     \For{$xag \in next\_xags$}
%%         \State $retQC$ + $parseXAG(xag, subQC)$
%%     \EndFor
%%     \State
%%     \Return $retQC$
%%     \EndFunction
%%   \end{algorithmic}
%% \end{algorithm}

\subsection{MCT Clean and Dirty Ancilla Decomposition Implications}

Observe that the XAG for the $n$-input AND function is the chain of AND nodes in~\figref{fig-mct-xag}.

\begin{figure}[t]
  \centering
  \scalebox{1.0} {
    \begin{tikzpicture}[
    gate/.style={circle, draw, inner sep=2pt},
    every node/.style={font=\small}
]

% Inputs
\node (x0) at (-0.5,0) {$x_0$};
\node (x1) at (0,0) {$x_1$};
\node (x2) at (0.5,0) {$x_2$};
\node at (1,0) {$\cdots$};
\node (xn) at (2,0) {$x_{n-1}$};

% AND nodes
\node[gate] (a1) at (0,0.5) {$\wedge$};
\node[gate] (a2) at (0.5,1) {$\wedge$};
\node[gate] (ak) at (2,1.5) {$\wedge$};

% Output
\node (out) at (2.5,2) {$x_0\cdots x_{n-1}$};

% Edges
\draw (x0) -- (a1);
\draw (x1) -- (a1);
\draw (x2) -- (a2);

\draw (a1) -- (a2);
\draw[dotted] (a2) -- (ak);

\draw (xn) -- (ak);

\draw[->] (ak) -- (out);

\end{tikzpicture}



  }
  \caption{XAG for an MCT gate}
  \label{fig-mct-xag}
\end{figure}

This allows us to create a a quantum circuit based off the gadgets in~\tabref{fig-gadgets}. This is depicted
in~\figref{fig-mct-gadgets}. 

\begin{figure}[t]
  \centering
  \scalebox{1.0} {
    \begin{quantikz}[row sep=0.2em, column sep=0.5em]
\lstick{$x_1$}                & \ctrl{1}                   & \ctrl{1}                           & \qw               \\
\lstick{$x_2$}                & \ctrl{2}                   & \ctrl{2}                           & \qw               \\
\setwiretype{n}               & \setwiretype{n}            & \setwiretype{n}                    & \setwiretype{n}   \\
\lstick{$x_3$}                & \gate[2]{\wedge}           & \gate[2]{\wedge^{\dagger}}         & \qw                 \\
\lstick{$a_1$}                & \ghost{\wedge}\ctrl{1}     & \ghost{\wedge^{\dagger}}\ctrl{1}   & \qw               \\
\setwiretype{n}               &                            &                                    & \setwiretype{n}   \\
\setwiretype{n}\push{\vdots}  & \push{\vdots}              & \push{\vdots}                      & \setwiretype{n}   \\
\setwiretype{n}               &                            &                                    & \setwiretype{n}   \\
\lstick{$x_n$}                & \gate[2]{\wedge}\ctrl{-1}] & \vqw{-1}\vqw{1}                    & \qw               \\
\lstick{$a_{n-2}$}            & \ghost{\wedge}\ctrl{1}     & \targ[style={double}]{}            & \qw               \\
\lstick{$f$}                  & \targ{}                    & \qw                                &
\end{quantikz}

%% \begin{quantikz}[row sep=0.3cm, column sep=0.55cm]
%% \qw              & \qw             & \qw        & \qw             & \qw              & \gate[2]{F_i}       & \qw                    & \qw       & \qw             & \qw      & \qw        & \qw              & \gate[2]{F_i}        & \qw                     & \qw        \\
%% \qw              & \qw             & \qw        & \qw             & \qw              & \ghost{F_i}         & \qw                    & \qw       & \qw             & \qw      & \qw        & \qw              & \ghost{F_i}          & \qw                     & \qw        \\
%% \qw              & \qw             & \qw        & \qw             & \ctrl{2}         & \qw                 & \ctrl{2}               & \qw       & \qw             & \qw      & \qw        & \ctrl{2}         & \qw                 & \ctrl{2}                & \qw         \\
%% \qw              & \qw             & \octrl{2}  & \qw             & \qw              & \qw                 & \qw                    & \qw       & \qw             & \qw      & \qw        & \qw              & \qw                 & \qw                     & \qw         \\
%% \lstick{$a_1$}   & \qw             & \qw        & \qw             & \gate{i\omega Z} & \targ{}             & \gate{i\omega Z^\dagger} & \qw       & \octrl{1}       & \qw      & \qw        & \gate{i\omega Z} & \targ{}             & \gate{i\omega Z^\dagger}  & \qw          \\
%% \lstick{$a_2$}   & \qw             & \gate{iZ}  & \gate{H}        & \ctrl{-1}        & \qw                 & \ctrl{-1}              & \gate{H}  & \gate{iZ^\dagger} & \qw      & \gate{H}   & \ctrl{-1}        & \qw                 & \ctrl{-1}               & \gate{H}     \\
%% \lstick{}        & \gate{H}        & \ctrl{-1}  & \qw             & \qw              & \qw                 & \qw                    & \qw       & \ctrl{-1}       & \gate{H} & \qw        & \qw              & \qw                 & \qw                     & \qw          \\
%%\end{quantikz}


  }
  \caption{Synthesis of~\figref{fig-mct-xag} using dirty ancilla gadgets}
  \label{fig-mct-gadgets}
\end{figure}

This construction is remarkably similar to the ``V-chain'' that is created when the clean ancilla 
decomposition is used, such as in~\figref{fig-clean-mct}. This demonstrates that the difference in
cost between a dirty ancilla decomposition and a clean ancilla decomposition is the cost of the
final $\wedge$ gadget, which uses one more T-gate than the Toffoli gate that the clean ancilla
decomposition uses.






%\input{Exp}

\vspace{0.5cm}
\section{Conc}
\label{Conc}

In this work, we have proposed a method of generating quantum circuits from XAG
using dirty ancilla, instead of the clean ancilla utilized by contemporary
methods. Using our specially proposed gadgets, we demonstrate that it differs
in cost from the clean ancilla methods by a single T-gate. We used it to demonstrate
that the clean and dirty ancilla decompositions of the MCT gate differ by a
single T-gate.

Some future work to consider is how to schedule its ancilla, just like the clean ancilla
method uses the pebbling game to conserve clean ancilla qubits.


% DCTODO: Add Evaluation, Results, and Conclusion Sections
\vspace{2cm}
\renewcommand{\baselinestretch}{1.10}
\bibliographystyle{unsrt}
\bibliography{ref}

\end{document}
